% Part: computability
% Chapter: computability-theory
% Section: equiv-ce-defs

\documentclass[../../../include/open-logic-section]{subfiles}

\begin{document}

\olfileid{cmp}{thy}{eqc}

\olsection[تعریف‌های مجموعه‌های م.ب.]{تعریف‌های هم‌ارز مجموعه‌های
  محاسبه‌پذیرِ برشمردنی}

قضیهٔ زیر چند بیان مهم و هم‌ارز از معنای محاسبه‌پذیرِ برشمردنی‌بودن
به دست می‌دهد.

\begin{thm}
\ollabel{thm:ce-equiv}
بگذارید $S$ مجموعه‌ای از اعداد طبیعی باشد. آنگاه گزاره‌های زیر
هم‌ارزند:
\begin{enumerate}
\item\ollabel{case:ce} $S$ محاسبه‌پذیرِ برشمردنی است.
\item\ollabel{case:ran-pc} $S$ برد یک تابع محاسبه‌پذیر \emph{جزئی} است.
\item\ollabel{case:ran-prim} $S$ تهی است یا برد یک تابع بازگشتی اولیه است.
\item\ollabel{case:ce-domain} $S$ \emph{دامنهٔ} یک تابع محاسبه‌پذیر جزئی است.
\end{enumerate}
\end{thm}

\begin{explain}
سه بند نخست می‌گویند که به‌طور هم‌ارز می‌توان هر مجموعهٔ ناتهی و
محاسبه‌پذیرِ برشمردنی را با تابعی محاسبه‌پذیر، تابعی محاسبه‌پذیر جزئی
یا تابعی بازگشتی اولیه برشمرد. بند چهارم می‌گوید اگر $S$
محاسبه‌پذیرِ برشمردنی باشد، برای نمایه‌ای چون~$e$ داریم
\[
S = \Setabs{x}{\cfind{e}(x) \fdefined}.
\]
به بیان دیگر، $S$ مجموعهٔ ورودی‌هایی است که محاسبهٔ $\cfind{e}$
برای آن‌ها متوقف می‌شود. به همین دلیل، مجموعه‌های محاسبه‌پذیرِ
برشمردنی را گاه \emph{نیمه‌تصمیم‌پذیر} می‌نامند: اگر عددی در مجموعه
باشد سرانجام پاسخ «آری» می‌گیرید، اما اگر نباشد هرگز پاسخ «نه»
نمی‌گیرید!{}
\end{explain}

\begin{proof}
چون هر تابع بازگشتی اولیه محاسبه‌پذیر و هر تابع محاسبه‌پذیر،
محاسبه‌پذیر جزئی است، \olref{case:ran-prim} مستلزم
\olref{case:ce} و \olref{case:ce} مستلزم~\olref{case:ran-pc} است.
(توجه کنید اگر $S$ تهی باشد، $S$ برد تابع محاسبه‌پذیر جزئی‌ای است که
هیچ‌جا تعریف نشده است.) اگر نشان دهیم \olref{case:ran-pc} مستلزم
\olref{case:ran-prim} است، هم‌ارزی سه بند نخست را نشان داده‌ایم.

پس فرض کنید $S$ برد تابع محاسبه‌پذیر جزئیِ~$\cfind{e}$ باشد. اگر $S$
تهی است، کار تمام است. وگرنه بگذارید $a$ عنصری دلخواه از~$S$ باشد.
بنا بر قضیهٔ صورت نرمال کلینی می‌توانیم بنویسیم:
\[
\cfind{e}(x) = U(\umin{s}{T(e, x, s)}).
\]
به‌ویژه، $\cfind{e}(x)\fdefined$ و برابر $y$ است اگر و تنها اگر
$s$ای وجود داشته باشد که $T(e,x,s)$ و $U(s)=y$. تابع $f(z)$ را چنین
تعریف کنید:
\[
f(z) = \begin{cases}
  U((z)_1) & \text{اگر $T(e, (z)_0, (z)_1)$} \\
  a        & \text{در غیر این صورت.}
\end{cases}
\]
آنگاه $f$ بازگشتی اولیه است، زیرا $T$ و~$U$ چنین‌اند.
\iftag{TMs}{برحسب ماشین‌های تورینگ، اگر $z$ زوج
  $\tuple{(z)_0,(z)_1}$ را چنان کد کند که $(z)_1$ محاسبهٔ متوقف‌شوندهٔ
  ماشین~$M_e$ روی ورودی $(z)_0$ باشد، آنگاه $f$ خروجی محاسبه را
  بازمی‌گرداند؛ وگرنه~$a$ را بازمی‌گرداند.}

باید نشان دهیم $S$ برد~$f$ است؛ یعنی برای هر عدد طبیعی~$y$،
$y\in S$ اگر و تنها اگر در برد~$f$ باشد. در جهت پیشرو فرض کنید
$y\in S$. پس $y$ در برد $\cfind{e}$ است و برای بعضی $x$ و~$s$،
$T(e,x,s)$ برقرار و $U(s)=y$. اما در این صورت
$y=f(\tuple{x,s})$. برعکس، فرض کنید $y$ در برد~$f$ باشد. آنگاه یا
$y=a$، یا برای یک~$z$، $T(e,(z)_0,(z)_1)$ و $U((z)_1)=y$. در حالت
دوم $\cfind{e}((z)_0)\fdefined=y$؛ پس در هر دو حالت $y\in S$.

(نماد $\cfind{e}(x)\fdefined=y$ یعنی «$\cfind{e}(x)$ تعریف شده و
برابر $y$ است.» می‌توانستیم به همان خوبی از
$\cfind{e}(x)=y$ استفاده کنیم، اما پیکان اضافی گاه یادآوری می‌کند که
با تابعی جزئی سروکار داریم.)

برای پایان برهان \olref{thm:ce-equiv} کافی است نشان دهیم
\olref{case:ce} و~\olref{case:ce-domain} هم‌ارزند. نخست نشان دهیم
\olref{case:ce} مستلزم~\olref{case:ce-domain} است. فرض کنید $S$ برد
تابع محاسبه‌پذیر~$f$ باشد، یعنی
\[
S = \Setabs{y}{\text{برای یک $x$، } f(x) = y}.
\]
بگذارید
\[
g(y) = \umin{x}{(f(x) = y)}.
\]
آنگاه $g$ تابعی محاسبه‌پذیر جزئی است و $g(y)$ تعریف شده است اگر و
تنها اگر برای یک~$x$، $f(x)=y$. به بیان دیگر، دامنهٔ~$g$ همان
برد~$f$ است. \iftag{TMs}{برحسب ماشین‌های تورینگ: با داشتن ماشین
تورینگ~$F$ که عناصر~$S$ را برمی‌شمارد، بگذارید $G$ ماشین تورینگی باشد
که با جست‌وجوی خروجی‌های~$F$ برای یافتن عنصر داده‌شده، $S$ را
نیمه‌تصمیم می‌گیرد؛ اگر عنصر را بیابد متوقف می‌شود و اگر نیابد تا ابد
به جست‌وجو ادامه می‌دهد.}{}

سرانجام، برای نشان‌دادن اینکه \olref{case:ce-domain} مستلزم
\olref{case:ce} است، فرض کنید $S$ دامنهٔ تابع محاسبه‌پذیر جزئیِ
$\cfind{e}$ باشد، یعنی
\[
S = \Setabs{x}{\cfind{e}(x) \fdefined}.
\]
اگر $S$ تهی است کار تمام است؛ وگرنه بگذارید $a$ عنصری دلخواه از~$S$
باشد. $f$ را چنین تعریف کنید:
\[
f(z) = \begin{cases}
(z)_0 & \text{اگر $T(e,(z)_0,(z)_1)$} \\
a & \text{در غیر این صورت.}
\end{cases}
\]
آنگاه مانند بالا، عدد $x$ در برد~$f$ است اگر و تنها اگر
$\cfind{e}(x)\fdefined$، یعنی اگر و تنها اگر $x\in S$.
\iftag{TMs}{برحسب ماشین‌های تورینگ: با داشتن ماشین~$M_e$ که $S$ را
نیمه‌تصمیم می‌گیرد، همهٔ محاسبه‌های ممکن ماشین تورینگ را مرور کنید و
ورودی‌های متناظر با محاسبه‌های متوقف‌شونده را بازگردانید تا عناصر $S$
برشمرده شوند.}{}
\end{proof}

بند~\olref{case:ce-domain} از \olref{thm:ce-equiv} راهی مناسب برای
برشمردن مجموعه‌های محاسبه‌پذیرِ برشمردنی فراهم می‌کند: به ازای هر~$e$،
بگذارید $W_e$ دامنهٔ $\cfind{e}$ باشد، یعنی
\[
W_e = \Setabs{x}{\cfind{e}(x) \fdefined}.
\]
آنگاه اگر $A$ هر مجموعهٔ محاسبه‌پذیرِ برشمردنی باشد، برای یک~$e$،
$A=W_e$.

قضیهٔ زیر توصیف دیگری از مجموعه‌های محاسبه‌پذیرِ برشمردنی به دست
می‌دهد.

\begin{thm}
\ollabel{thm:exists-char}
مجموعهٔ $S$ محاسبه‌پذیرِ برشمردنی است اگر و تنها اگر رابطه‌ای
محاسبه‌پذیر چون $R(x,y)$ وجود داشته باشد که
\[
S = \Setabs{ x }{ \lexists[y][R(x,y)] }.
\]
\end{thm}

\begin{proof}
در جهت پیشرو فرض کنید $S$ محاسبه‌پذیرِ برشمردنی باشد. آنگاه برای یک
$e$، $S=W_e$. برای این مقدار~$e$ می‌توان نوشت:
\[
S = \Setabs{ x }{ \lexists[y][T(e, x, y)] }.
\]
در جهت عکس، فرض کنید
$S=\Setabs{x}{\lexists[y][R(x,y)]}$. $f$ را چنین تعریف کنید:
\[
f(x) \simeq \umin{y}{R(x, y)}.
\]
آنگاه $f$ محاسبه‌پذیر جزئی و $S$ دامنهٔ~$f$ است.
\end{proof}

\end{document}
